\scanpage{74}
\uline{Théorème 1}$_n$ Soit $A$ anneau str. local noeth., de dim $n$, \inserted{\deleted{\unsure{str. excellent}} satisfaisant cond. de \unsure{finitude|stricte}}
\deleted{alors pour tout} \inserted{$Y=\operatorname{Spec}(A)$} $F$ un faisceau de $\ell$ torsion sur
$X=\operatorname{Spec} A[t]$, \inserted{$\ell$ premier à la car. résiduelle,} tel que $x\in\operatorname{Supp} F\Longrightarrow\dim\underline O_{X,x}\geq m$.
(en un mot ; $F$ est ``nul en codim $<m$). \deleted{Dans
ce cas \unsure{si|où} \unsure{moins|montre} \illegible{1} \unsure{que} en plus \unsure{que} support \unsure{en codim} \unsure{pour} \unsure{le pt gén.} de $X$}.
Alors $H^i(X,F)=0$ pour $i>n+1-m$.
% slanted
\margin{\uline{Question} Est-il \unsure{nécessaire} que $F$ soit de torsion ???}
[En particulier, $\operatorname{cd}_{\ell}(X)\leq n+1$\deleted{\unsure{]}} , et pour tte partie fermée $Z$
de $X$, \deleted{de codim} $\operatorname{cd}_{\ell}(Z)\leq n+1-\operatorname{codim}(Z,X)$]

\uline{Théorème} $2_n$ Soit $A$ anneau str. local noeth de
dim $n$, \deleted{et $f\in A$} $X=\operatorname{Spec}(A)$, $f\in A$, $U=X_f$, Alors
$\operatorname{cd}_{\ell}(U)\leq n$ pour tout $\ell$ premier à la car. résiduelle.

[N.B. $1_n\Rightarrow2_n$ car on plonge $U$ comme fermé \inserted{de codim $\geq1$,} dans $Y[t]$.]

\uline{Lemme 1} $\underbrace{2_{n'}\ \text{pour }n'\leq n+1}_{2'_{n+1}}\Longrightarrow 1_n$
\inserted{(aussi $=1_{n'}$ \unsure{pour} $n'\leq n$ i.e. $1'_n$)}. En effet, on est ramené à
prouver que si $S\subset X$ fermé est de codim $m$,
alors $\operatorname{cd}_{\ell}(S)\leq n+1-m$. \inserted{On \unsure{peut supposer} $S$ irréd.} Or soit $\bar S$ l'adhérence de
$S$ dans $\bar X=\mathrm P^1_Y$, considérons $T=\bar S-S$, c'est un
schéma str. local ($\subset\bar X-X$), soit $x$ son pt fermé.
Cons. \deleted{pour} \inserted{\unsure{tout|pour}} $F$ de $\ell$-torsion sur $S$
\[
H^i(\bar S,F)\longrightarrow H^i(S,F)\longrightarrow H_T^{i+1}(\bar S,F)
\]
Or
\[
\begin{aligned}
H^i(\bar S,F)&=H^i(\bar S_0,F)=H^i(\mathrm P^1_k,F_0),\ (\text{nul si }i\geq3)\\
H_T^{i+1}(\bar S,F)&=H^i(\bar S_x-T_x,F_{(x)})\qquad\text{si } i\geq2
\end{aligned}
\]
\deleted{nul par l'hyp. \unsure{de}}
D'ailleurs $\dim\bar S_x\leq n+1-m$ car \deleted{\unsure{rien|bien}}
% archive sheet number: 73
